\subsection{Countermodule}

DEFINITION 3. --- Let A be a ring, M an A-module, and \(E = \text{End}_{A}(M)\) the endomorphism ring of M. The countermodule of M is the left E-module with the same underlying additive group as M and external law \((c, x) \mapsto c(x)\) .

Let Z be the center of the ring A. For every \(a \in Z\) , the homothety \(a_{M}\) belongs to E. Consequently, E is canonically endowed with the structure of a Z-algebra. In particular, if M is a finitely generated Z-module, then the countermodule of M is finitely generated.

Lemma 4. --- Let M be a left A-module with finitely generated countermodule. There exist a natural number m and an injective \(A_{M}\) -linear mapping from \((\mathrm{A}_{\mathrm{M}})_{s}\) to \(M^{m}\) .

Set \(\mathrm{E}=\mathrm{End}_{\mathrm{A}}(\mathrm{M})\) . Let \((x_{1},\ldots,x_{m})\) be a finite generating family of the E-module M. The mapping \(\varphi:a\mapsto(ax_{1},\ldots,ax_{m})\) from \((\mathrm{A}_{\mathrm{M}})_{s}\) to \(M^{m}\) is \(A_{M}\) -linear. Let a be an element of \(A_{M}\) such that \(\varphi(a)=0\) . The set of elements x of M such that ax = 0 is an E-submodule of M containing \(x_{1}, \ldots, x_{m}\) and is therefore equal to M, which implies a = 0.

PROPOSITION 6. --- Let M be an Artinian (resp. Noetherian) left A-module with finitely generated countermodule. The ring of homotheties \(A_{M}\) of M is left Artinian (resp. left Noetherian).

This follows from Lemma 4 and Proposition 3 of VIII, p.~3.

COROLLARY. --- Let A be a commutative ring.

\begin{enumerate}
\def\labelenumi{\alph{enumi})}
\item
  Let M be a Noetherian A-module. The ring \(A_{M}\) is Noetherian.
\item
  Let M be an A-module of finite length. The ring \(A_{M}\) is Artinian.
\end{enumerate}

Let M be an A-module. Under the assumptions of a) or b), the A-module M is finitely generated. Since A is commutative, \(A_{M}\) is contained in the ring \(\operatorname{End}_{\mathrm{A}}(\mathrm{M})\) , so that the countermodule of M is finitely generated. It then suffices to apply Proposition 6.

Remark. --- Let A be a ring. An Artinian left A-module M with finitely generated countermodule has finite length: indeed, the ring of homotheties \(A_{M}\) of M is left Artinian (Proposition 6), and M is an Artinian module over \(A_{M}\) ; by VIII, p.~7, Proposition 4, the module M has finite length over \(A_{M}\) and therefore also over A.

In particular, every finitely generated Artinian module over a commutative ring has finite length. By contrast, a finitely generated Artinian module over a noncommutative ring is not necessarily of finite length (VIII, p.~16, Exercise 12).

